%! Diagonalizing a Matrix — Unit 6, Lesson 6.2 — Student Packet Cover Sheet
\documentclass[10pt]{article}
\usepackage{linalg-article}
\usepackage{linalg-boxes}
\usepackage{ltablex}
\keepXColumns

\begin{document}

% Full-bleed burgundy banner
\begin{tikzpicture}[remember picture, overlay]
  \fill[burgundy] (current page.north west)
    rectangle ([yshift=-0.9in]current page.north east);
\end{tikzpicture}
\vspace{-0.8in}

\begin{center}
  {\color{white}\textbf{\LARGE Linear Algebra}}\\[4pt]
  {\color{blushmid}\large\textbf{Unit 6 \quad Eigenvalues and Eigenvectors}}\\[6pt]
  {\color{white}\textbf{\large Lesson 6.2 \quad Diagonalizing a Matrix}}
\end{center}

\vspace{0.22in}
\namedateperiod
\vspace{0.18in}

\begin{learningtargetbox}
\small
\begin{enumerate}[leftmargin=1.5em, itemsep=5pt]
  \item \ldots{} line a matrix's eigenvectors into the columns of $X$ and its eigenvalues down the diagonal of $\Lambda$, and explain why $AX=X\Lambda$.
  \item \ldots{} \textbf{diagonalize} a matrix as $A=X\Lambda X^{-1}$ (when it has enough independent eigenvectors).
  \item \ldots{} use $A^k=X\Lambda^k X^{-1}$ to compute powers of a matrix quickly, and say when a matrix \emph{cannot} be diagonalized.
\end{enumerate}
\end{learningtargetbox}

\vspace{0.14in}

\begin{tocbox}
\small
\renewcommand{\arraystretch}{1.5}
\begin{tabularx}{\linewidth}{@{} c l X r @{}}
\rowcolor{blush}
\textbf{\#} & \textbf{Component} & \textbf{Description} & \textbf{Score} \\
\midrule
1 & Warm-Up        & Spiral review: build a matrix from columns, the $2\times2$ inverse, scaling columns & \blank{1.2cm} \\
2 & Guided Notes   & $AX=X\Lambda$; the factorization $A=X\Lambda X^{-1}$; powers $A^k=X\Lambda^k X^{-1}$ & \blank{1.2cm} \\
3 & Group Activity & Building $X,\Lambda$; a full diagonalization + a power; a matrix that \emph{can't} be diagonalized & \blank{1.2cm} \\
4 & Exit Ticket    & Quick check: diagonalize a matrix and explain why powers become easy & \blank{1.2cm} \\
5 & Homework       & Diagonalize, compute a power, decide diagonalizability, justify the method & \blank{1.2cm} \\
\midrule
  & \multicolumn{2}{r}{\textbf{Total}} & \blank{1.2cm} \\
\end{tabularx}
\end{tocbox}

\vspace{0.10in}

\begin{remindbox}
\small Yesterday you learned to \emph{find} a matrix's eigenvalues and eigenvectors. Today you put them
to work. Line the eigenvectors up as the columns of a matrix $X$ and the eigenvalues down the diagonal
of $\Lambda$; then every matrix with enough independent eigenvectors factors as
$A=X\Lambda X^{-1}$ --- it \emph{acts like the diagonal matrix} $\Lambda$ in eigenvector coordinates. The
payoff: powers collapse to $A^k=X\Lambda^k X^{-1}$, so raising $A$ to a power just means raising each
eigenvalue to that power.
\end{remindbox}

\end{document}
